\section{Future work}

To address identified limitations, we propose the following enhancements.

\textbf{Thesis Vertex Labeling}. Each thesis vertex will receive dual categorization via two distinct labels: \textit{Episode} and \textit{Temporal}. The \textit{Episode} labeling classifies thesis formulations: FACT -  factual claims verifiable through independent evidence; OPINION - subjective opinions, requiring contextual interpretation;  PREDICTION - speculative predictions lacking immediate verification. Meanwhile, \textit{Temporal} labeling specifies the duration, over which a statement remains relevant: STATIC - statically enduring facts; DYNAMIC - temporally limited assertions expiring upon subsequent developments; ATEMPORAL - universally applicable truths unaffected by chronology.

\textbf{Time Interval Specification}. We introduce explicit timestamps for each thesis vertex, identifying its creation (t\_created), validity onset (t\_valid), expiration (t\_expired), invalidation (t\_invalid) and potential override by newer information (invalidated\_by). These parameters facilitate precise tracking of knowledge lifecycle stages~\cite{Rasmussen2025ZepAT}.

\textbf{Fixed Predicate Fields}. Text fields in simple triples predicates adopt predefined, time-independent values from verified and periodically updated collection (glossary). Object vertices store additional metadata such as entity types and brief descriptions, while predicates include textual representations specifying relationships between subjects and objects~\cite{zhang-soh-2024-extract}.

These modifications collectively aim to enhance reliability, scalability, and usability of the proposed method, thereby mitigating current drawbacks effectively.
